\chapter{Applications to Fusion, Plasma Radiation Sources, and Biological Exposure}
\label{ch:applications}

\section{Proton--boron-11 fusion}

The \pB \ reaction
\begin{equation}
p+{}^{11}\mathrm{B}\rightarrow 3\alpha+8.68\ \mathrm{MeV}
\end{equation}
is attractive because the primary products are charged and the fuel constituents are comparatively abundant. The reaction is often described as aneutronic, but a complete reactor assessment must include secondary reactions, impurities, activation, and high-energy photons. The central energy-balance obstacle is that useful reactivity requires temperatures at which electron radiation is severe \cite{rider1995,nevins2000,sikora2016,ochs2022}.

Distribution engineering addresses the problem at its physical origin. A reactor may seek hot or fast ions near the reaction-energy range while maintaining cooler electrons, redirecting alpha energy into ions, and suppressing the superthermal electron tail. Magnetic perturbations could contribute by selectively deconfining radiatively costly electrons or by coupling them to wave and electrostatic structures that transfer energy elsewhere. However, electrons also provide ambipolar confinement, current, quasi-neutrality, and collisional coupling. Overaggressive removal can deepen potentials, drive replacement heating, destabilize the plasma, or increase ion losses.

The radiobiological extension matters for plant design. Even if \pB \ reduces neutron shielding requirements relative to deuterium--tritium, bremsstrahlung can dominate prompt dose outside the plasma. Source-term control may reduce shield thickness, maintenance dose, diagnostic saturation, and occupational exposure. These benefits must be quantified spectrally because low-energy X rays are easily shielded but can produce high local dose near thin windows, while hard photons dominate deep penetration and streaming.

\section{Mirror and open-field-line systems}

Mirror machines provide a natural environment for energy-selective escape. The loss cone, mirror ratio, ambipolar potential, and relativistic confinement scaling determine which electrons leave \cite{pastukhov1974,ochs2023mirror}. A controlled perturbation can enlarge the loss region for a selected phase-space population. The engineering advantage is that open ends provide intentional collection geometry; the disadvantage is that end losses can be large and produce intense localized radiation when electrons strike hardware.

A mirror implementation should include an expanding magnetic nozzle, electrostatic energy recovery where feasible, and low-$Z$ electron collectors. Diagnostics placed along the end loss channel can directly measure the removed distribution, providing a strong validation of the cutoff mechanism. Biological dose should be evaluated both radially and along the axis, where streaming is most likely.

\section{Tokamaks and stellarators}

Resonant magnetic perturbations are established tools for modifying edge topology in toroidal confinement \cite{evans2004,fitzpatrick1993,wesson2011}. Applying them to bremsstrahlung suppression is more difficult because the energetic population may be core localized, while broad stochasticity degrades energy confinement. A feasible concept would target an energetic-electron subpopulation with mode structure, frequency, and radial localization that minimize bulk transport.

Runaway-electron physics provides a cautionary analogy. Synchrotron and bremsstrahlung reaction can limit runaway energy, while magnetic perturbations can deconfine runaways, but localized wall impacts create severe material loads. The same trade-off appears here: removal of radiating electrons must not convert a diffuse photon source into an unacceptable localized electron beam. Bremsstrahlung reaction also reshapes the distribution and should be retained in any runaway-like regime \cite{embreus2016}.

\section{Laser-produced plasmas}

High-intensity lasers generate nonthermal electrons through resonance absorption, $\mathbf{J}\times\mathbf{B}$ heating, wakefield processes, and sheath acceleration. These electrons create bremsstrahlung when traversing dense plasma or solid targets. Laser-driven \pB \  experiments have demonstrated energetic alpha production, but conversion efficiency, electron heating, target geometry, and radiation losses remain central challenges \cite{labaune2013,picciotto2014,giuffrida2020}.

External magnetic fields can guide electron transport, modify filamentation, and alter refluxing in the target. The proposed framework can evaluate whether a magnetic field reduces source bremsstrahlung in the interaction region or merely redirects electrons to a new converter location. Because laser sources are strongly pulsed and directional, time-resolved dosimetry and angular spectra are indispensable. Extremely high instantaneous dose rate does not automatically imply a beneficial FLASH response in nearby tissue; the field may be spatially nonuniform, mixed with particles, and delivered at a total dose outside therapeutic conditions.

\section{Dense plasma focus and Z-pinch devices}

Dense plasma focus and Z-pinch systems produce intense pulsed X rays, electron beams, ions, and in some operating modes neutrons \cite{haines2011,lee2014}. Their compactness has motivated study as radiography and ultra-high-dose-rate radiation sources \cite{sumini2019,isolan2022}. Magnetic perturbations are inherent in the pinch dynamics, and instability controls may alter both plasma confinement and radiation output.

For radiobiology, these devices are valuable test beds because they naturally couple plasma physics to pulsed exposure. They also illustrate the need for mixed-field characterization. A measured biological effect cannot be attributed to X rays without excluding electrons, ions, neutrons, ultraviolet radiation, electromagnetic pulse artifacts, sample heating, and mechanical shock. The full methodology in \Cref{ch:methodology} is especially relevant.

\section{Accelerator and beam-plasma systems}

Relativistic electron beams interacting with plasma or high-$Z$ targets generate tunable bremsstrahlung. Wakefield accelerators and beam-driven plasma sources can produce compact ultrashort X-ray pulses. Magnetic optics shape the electron beam before conversion; plasma-based magnetic structures may also influence energy spread and divergence. Distribution-function engineering can reduce unwanted halo electrons or select the energy range that contributes to a desired photon spectrum.

In an intentional radiation source, ``suppression'' is not always the objective. The same framework can optimize spectral quality: reduce biologically inefficient or shielding-intensive components while retaining photons useful for imaging, therapy, or materials testing. The biological weighting in Equation~\eqref{eq:bio-objective} can be replaced by image-quality or detector-response weighting when appropriate.

\section{Medical and preclinical radiation research}

Plasma-generated photon sources may serve radiobiology or radiotherapy research if they can deliver calibrated, reproducible fields. Their possible advantages include compact geometry, short pulses, and unconventional spectra. Their barriers include source instability, dose-rate metrology, field uniformity, mixed radiation, and regulatory complexity.

A source designed for preclinical use should be benchmarked against conventional kilovoltage or megavoltage X rays using traceable dosimetry and biological endpoints. The linear-quadratic model, oxygen effect, Relative Biological Effectiveness (RBE), and tissue constraints remain relevant, but nonstandard pulse structure may require additional modeling. A direct transition to patient treatment would be premature without extensive dosimetric, toxicological, engineering, and clinical validation.

\section{Space and aerospace applications}

High-energy plasma and accelerator systems used in space propulsion, radiation testing, or compact power concepts can create ionizing radiation. Space radiation biology is usually dominated by trapped particles, solar particles, and galactic cosmic rays, but local device-generated photons and particles can add to crew or electronics dose. The source-to-biology framework supports onboard optimization: magnetic topology can be evaluated simultaneously for device performance, vehicle shielding, crew location, and mission duty cycle.

For crew protection, equivalent and effective dose are more relevant than in vitro survival. Mixed fields require radiation-quality weighting and organ geometry. The framework also applies to ground testing, where pulsed plasma devices may expose operators through direct lines of sight or shield penetrations.

\section{Occupational radiation protection}

A facility safety analysis begins with source identification and proceeds through engineered controls, administrative controls, and personal monitoring. The protection principles are justification, optimization, and limitation, implemented through time, distance, shielding, interlocks, access control, surveys, and ALARA planning \cite{icrp103,iaeaGSR3}. Magnetic source suppression is an engineered control, but it does not eliminate the need for passive shielding or monitoring because the control can fail or operate outside its validated range.

Operational dose should be evaluated for normal operation, commissioning, maintenance, credible faults, and post-shutdown activation. Source suppression may reduce prompt dose but not residual activation. Conversely, reduced photon heating can lower activation indirectly by permitting different materials or operating conditions. These couplings are facility specific.

\section{Environmental and public exposure}

Public exposure from a well-designed plasma facility should be dominated by controlled emissions and shielding leakage rather than direct access to the source. The relevant analysis uses annual operation, occupancy, skyshine where applicable, effluents, activation, and accident scenarios. A reduction in bremsstrahlung source term can reduce off-site dose, but the public-dose consequence must be calculated with site geometry and regulatory models rather than inferred from laboratory detector ratios.

Environmental radiobiology also considers non-human biota, radionuclide transport, and ecosystem effects \cite{baatout2023,unscear2020}. For a nominally aneutronic system, radionuclide inventories may be smaller than in neutron-rich fusion, but secondary sources and activated components still require characterization.

\section{Electronics and instrumentation}

Radiation damage is not only biological. X rays and secondary electrons cause ionization in detectors, cables, insulators, semiconductors, and control electronics. Total ionizing dose, single-event effects from secondary particles, displacement damage, and detector saturation can limit operation. A spectral source reduction may extend component life and improve diagnostic dynamic range even when worker dose is already acceptable.

The same transport matrix can produce component dose. Biological weighting is replaced by material energy deposition, displacement-damage functions, or device response. This dual use strengthens the case for preserving spectral information throughout the model.

\section{Design case study: source suppression versus additional shielding}

Consider two options for reducing dose at a maintenance location:
\begin{description}
\item[Option A:] magnetic control reduces the hard-photon source by a factor $R_S(E)$ but adds coil power, complexity, and electron collector load;
\item[Option B:] additional passive shielding changes the transport operator but leaves plasma performance unchanged.
\end{description}
The decision compares lifecycle cost, reliability, dose, mass, access, heat, activation, and maintainability. Source suppression may be superior for streaming through many penetrations because it acts before geometric branching. Shielding may be superior if the perturbation degrades fusion performance or creates new local sources. A hybrid design often dominates: moderate source reduction followed by optimized graded shielding.

\section{Summary}

The unified framework applies across advanced fusion, mirror systems, toroidal devices, laser plasmas, dense plasma focus sources, accelerator systems, medicine, space, and occupational protection. The specific objective changes, but the logic remains constant: control the particle distribution, calculate the complete external source, transport it to a defined receptor, and use an endpoint-appropriate response model.
